\subsection{How Do I Protect Early Gains From Collapse?}

If you have just started moving in the right direction, the main job is not to
push harder in a heroic way. The main job is to keep the new target state from
losing its early foothold while the old framework inertia is still nearby.
Early gains usually collapse when the barrier is allowed to rise again before a
simple repeatable routine has been installed. Protect the gain by lowering the
cost of the next action, stabilizing the environment, and making the next
decision window easier to enter than the old one.

\begin{enumerate}[nosep]
  \item Name the exact early gain you want to preserve. Do not protect a vague
        intention; protect a visible change such as ``I started studying at
        8:30'' or ``I walked after work twice this week.''
  \item Identify the first condition that would let the old pattern return.
        This is usually a return of clutter, ambiguity, delay, or fatigue, all
        of which raise the barrier.
  \item Install one cheap stabilizer before the next session. Put the book on
        the desk, set the clothes out, preload the file, or write the first
        line so the target state starts from a stronger base.
  \item Make the next action smaller than the previous one. Early gains are
        fragile because volitional energy is still limited; a smaller next step
        keeps the system moving without forcing a costly restart.
  \item Repeat the same entry rule at least several times in a row. The point
        is to make the new pattern feel ordinary before the old framework
        inertia has time to reclaim the space.
  \item Check the next decision window in advance. If the next opening is
        obvious, prepare for it before it arrives rather than trying to
        improvise after the barrier is already high again.
\end{enumerate}

\begin{remark}
Suppose you finally manage to study for twenty minutes after work. The gain is
not protected by congratulating yourself; it is protected by making tomorrow's
start cheaper than today's. If the laptop stays open, the notes stay visible,
and the first task is already written down, the next session begins with less
friction and the new behaviour has a better chance of surviving.
\end{remark}
